\subsection{Extension of scalars: case of commutative rings}

Now let A be a commutative ring, B a ring and \(\rho: \mathbf{A} \to \mathbf{B}\) a ring homomorphism such that \(\rho(\mathbf{A})\) is contained in the centre of B; in other words, \(\rho\) defines on B an

A-algebra structure. For every A-module E, the right B-module \(\mathrm{E}_{(\mathbf{B})} = \mathrm{E} \otimes_{\mathbf{A}} \mathbf{B}\) is then identified with \(B \otimes_{A} E\) , the A-module structures of \(\rho_{*}(B_{s})\) and \(\rho_{*}(B_{d})\) being identical by hypothesis. Recall that for every ordered pair (E, F) of A-modules, we have defined a canonical B-homomorphism

\[
\omega : \left(\operatorname{Hom} _ {\mathbf {A}} (\mathrm{E}, \mathrm{F})\right) _ {(\mathbf {B})} \rightarrow \operatorname{Hom} _ {\mathbf {B}} \left(\mathrm{E} _ {(\mathbf {B})}, \mathrm{F} _ {(\mathbf {B})}\right)\tag{10}
\]

such that for all \(u \in \operatorname{Hom}_{\mathbf{A}}(\mathrm{E}, \mathrm{F})\) , \(\omega(u \otimes 1) = u \otimes 1_{\mathbf{B}}\) (Algebra, Chapter II, § 5, no. 3).

PROPOSITION 11. Let A be a commutative ring, B a ring, \(\rho\) a homomorphism of A to the centre of B, and E and F two A-modules. Suppose that B is a flat A-module and E is finitely generated (resp. finitely presented). Then the canonical homomorphism (10) is injective (resp. bijective).

As \(\omega\) is composed of the canonical isomorphism

\[
\operatorname{Hom} _ {\mathbf {A}} (\mathrm{E}, \mathbf {B} \otimes_ {\mathbf {A}} \mathrm{F}) \rightarrow \operatorname{Hom} _ {\mathbf {B}} (\mathrm{E} _ {(\mathbf {B})}, \mathrm{F} _ {(\mathbf {B})})
\]

and the canonical homomorphism (8)

\[
\nu \colon \mathbf {B} \otimes_ {\mathbf {A}} \operatorname{Hom} _ {\mathbf {A}} (\mathbf {E}, \mathbf {F}) \rightarrow \operatorname{Hom} _ {\mathbf {A}} (\mathbf {E}, \mathbf {B} \otimes_ {\mathbf {A}} \mathbf {F})
\]

(loc. cit.), the proposition is a consequence of Proposition 10 of no. 9.

Suppose now that A and B are commutative and consider three A-modules \(E_{1}\) , \(E_{2}\) , \(E_{3}\) and an A-bilinear mapping \(f: E_{1} \times E_{2} \to E_{3}\) . Then there exists one and only one B-bilinear mapping \(f_{B}: E_{1(B)} \times E_{2(B)} \to E_{3(B)}\) such that

\[
f _ {\mathrm{B}} (1 \otimes x _ {1}, 1 \otimes x _ {2}) = 1 \otimes f (x _ {1}, x _ {2})
\]

for all \(x_{1} \in \mathbf{E}_{1}, x_{2} \in \mathbf{E}_{2}\) (Algebra, Chapter IX, § 1, no. 4, Proposition 1).

In the statement which follows we shall suppose that B is a flat A-module and, for every submodule \(E'\) of \(E_{i}\) (i = 1, 2, 3), we shall canonically identify \(E'_{(B)}\) with its image in \(E_{i(B)}\) (no. 3, Remark 2).

PROPOSITION 12. Let A, B be commutative rings, \(\rho\) a homomorphism of A to B, E \(_{1}\) , E \(_{2}\) , E \(_{3}\) three A-modules, \(f\) : E \(_{1} \times\) E \(_{2} \to\) E \(_{3}\) an A-bilinear mapping and

\[
f _ {\mathbf {B}} \colon \mathrm{E} _ {1 (\mathbf {B})} \times \mathrm{E} _ {2 (\mathbf {B})} \rightarrow \mathrm{E} _ {3 (\mathbf {B})}
\]

its extension. Consider a submodule \(\mathbf{F}_2\) of \(\mathbf{E}_2\) , a submodule \(\mathbf{F}_3\) of \(\mathbf{E}_3\) , and denote by \(\mathbf{T}\) the submodule of \(\mathbf{E}_1\) consisting of those \(x_1 \in \mathbf{E}_1\) such that \(f(x_1, x_2) \in \mathbf{F}_3\) for all \(x_2 \in \mathbf{F}_2\) . Suppose that \(\mathbf{B}\) is a flat A-module and that \(\mathbf{F}_2\) is finitely generated. Then \(\mathbf{T}_{(\mathbf{B})}\) is the set of those \(x_1' \in \mathbf{E}_{1(\mathbf{B})}\) such that \(f_{\mathbf{B}}(x_1', x_2') \in \mathbf{F}_{3(\mathbf{B})}\) for all \(x_2' \in \mathbf{F}_{2(\mathbf{B})}\) .

Let \(p\) be the canonical surjection \(\mathrm{E}_3 \to \mathrm{E}_3 / \mathrm{F}_3\) ; with each \(x_1 \in \mathrm{E}_1\) we associate the A-linear mapping \(x_2 \mapsto p(f(x_1, x_2))\) of \(\mathrm{F}_2\) to \(\mathrm{E}_3 / \mathrm{F}_3\) , which we denote by \(g(x_{1})\) ; then g is an A-homomorphism of \(E_{1}\) to \(\operatorname{Hom}_{\mathbf{A}}(\mathbf{F}_{2}, \mathbf{E}_{3}/\mathbf{F}_{3})\) and the kernel of g is precisely T. Since B is a flat A-module, we have the exact sequence

\[
0 \longrightarrow \mathrm{T} _ {(\mathbf {B})} \longrightarrow \mathrm{E} _ {1 (\mathbf {B})} \stackrel {1 \otimes g} {\longrightarrow} \left(\operatorname{Hom} _ {\mathbf {A}} \left(\mathrm{F} _ {2}, \mathrm{E} _ {3} / \mathrm{F} _ {3}\right)\right) _ {(\mathbf {B})}
\]

(no. 3, Proposition 1). By Proposition 11 the canonical homomorphism

\[
\omega : \left(\operatorname{Hom} _ {\mathbf {A}} (\mathrm{F} _ {2}, \mathrm{E} _ {3} / \mathrm{F} _ {3})\right) _ {(\mathbf {B})} \rightarrow \operatorname{Hom} _ {\mathbf {B}} (\mathrm{F} _ {2 (\mathbf {B})}, \left(\mathrm{E} _ {3} / \mathrm{F} _ {3}\right) _ {(\mathbf {B})})
\]

is injective. On the other hand, as B is a flat A-module, \((\mathrm{E}_{3}/\mathrm{F}_{3})_{(\mathrm{B})}\) is canonically identified with \(E_{3(\mathrm{B})}/F_{3(\mathrm{B})}\) ; taking the composition of \(\omega\) with \(1 \otimes g\) , we obtain a homomorphism u for which the sequence \(^{1}\)

\[
0 \longrightarrow \mathrm{T} _ {(\mathrm{B})} \longrightarrow \mathrm{E} _ {1 (\mathrm{B})} \stackrel {u} {\longrightarrow} \operatorname{Hom} _ {\mathrm{B}} \left(\mathrm{F} _ {2 (\mathrm{B})}, \mathrm{E} _ {3 (\mathrm{B})} / \mathrm{F} _ {3 (\mathrm{B})}\right)
\]

is exact. It follows immediately from the definitions that \(u(x_{1}^{\prime})\) , where

\[
x _ {1} ^ {\prime} = 1 \otimes x _ {1} \in \mathrm{E} _ {1 (\mathbf {B})},
\]

is the linear mapping which maps each \(x_2' \in \mathrm{F}_{2(\mathbf{B})}\) to the class mod. \(\mathrm{F}_{3(\mathbf{B})}\) of \(f_{\mathbf{B}}(x_1', x_2')\) ; by linearity this is also true for all \(x_1' \in \mathrm{E}_{1(\mathbf{B})}\) ; since the kernel of \(u\) is \(\mathrm{T}_{(\mathbf{B})}\) , the proposition is proved.

COROLLARY 1. Let A, B be two commutative rings, \(\rho: \mathbf{A} \to \mathbf{B}\) a homomorphism such that B is a flat A-module and E a finitely presented A-module. For every finitely generated submodule F of E, the submodule of the dual of \(\mathrm{E}_{(\mathrm{B})}\) orthogonal to \(\mathrm{F}_{(\mathrm{B})}\) is equal to \((\mathrm{F}')_{(\mathrm{B})}\) , where \(\mathrm{F}'\) is the submodule of the dual \(\mathrm{E}^*\) of E orthogonal to F.

By Proposition 11 \((\mathbf{E}^{*})_{(\mathbf{B})}\) is canonically isomorphic to the dual \((\mathbf{E}_{(\mathbf{B})})^{*}\) of \(\mathbf{E}_{(\mathbf{B})}\) . Then it suffices to apply Proposition 12 with \(\mathbf{E}_1 = \mathbf{E}^*\) , \(\mathbf{E}_2 = \mathbf{E}\) , \(\mathbf{E}_3 = \mathbf{A}\) , \(\mathbf{F}_2 = \mathbf{F}\) , \(\mathbf{F}_3 = \{0\}\) , and \(f\) the canonical bilinear form on \(\mathbf{E}^{*} \times \mathbf{E}\) .

COROLLARY 2. Let A, B be two commutative rings, \(\rho: A \to B\) a homomorphism such that B is a flat A-module and E an A-module. Then, for every finitely generated submodule F of E, the annihilator of \(F_{(B)}\) is the ideal aB of B, where a is the annihilator of F in A.

It suffices to apply Proposition 12 with \(\mathbf{E}_1 = \mathbf{A}\) , \(\mathbf{E}_2 = \mathbf{E}_3 = \mathbf{E}\) , \(\mathbf{F}_2 = \mathbf{F}\) , \(\mathbf{F}_3 = \{0\}\) .

Remark. If there is no ambiguity over the modules \(E_{t}\) nor the bilinear mapping f, \(F_{3}: F_{2}\) is sometimes used to denote the module denoted by T in Proposition 12 and it is called the transporter of \(F_{2}\) to \(F_{3}\) . The conclusion of Proposition 12 then reads

\[
\mathrm{F} _ {3 (\mathrm{B})}: \mathrm{F} _ {2 (\mathrm{B})} = \left(\mathrm{F} _ {3}: \mathrm{F} _ {2}\right) _ {(\mathrm{B})}.\tag{11}
\]

In the particular case when the \(E_{i}\) are equal to the ring A, f is multiplication and the \(F_{i}\) ideals \(a_{i}\) , we obtain the transporter formula

\[
\mathbf {B} \left(a _ {3}: a _ {2}\right) = \mathbf {B} a _ {3}: \mathbf {B} a _ {2}\tag{12}
\]

valid when B is a flat A-module and \(a_{2}\) is a finitely generated ideal.
% semantic-fragments: legacy-input-seam
\relax{}
